\section{Methodology}

\subsection{Problem Setup}

Consider a classification task where inputs $x$ have class label $y$ and group membership $g$. Groups arise from the interaction of core features (predictive of $y$) and spurious features (correlated with $y$ in training but not causally related). Under ERM training, models exploit spurious correlations because they provide simpler, more consistent signals.

\subsection{Loss Trajectory Detection}

For each sample $i$, we track loss $L_i(t)$ across training epochs $t$. Rather than computing onset delay (which requires tracking from epoch 0), we use loss at early detection epoch $T_{\text{early}} = 5$ (5\% of 100 epochs).

We identify candidate minority samples as those with loss above the $k$-th percentile:
\begin{equation}
\text{minority\_candidate}_i = \mathbf{1}[L_i(T_{\text{early}}) > \text{percentile}_k(L(T_{\text{early}}))]
\end{equation}

At 5\% minority rate, $k = 95$ selects the top 5\% of high-loss samples.

\subsection{Linear Probe Analysis}

To verify the simplicity bias mechanism, we train linear probes on frozen ResNet-18 features at multiple epochs. Two separate logistic regression probes are trained:
\begin{itemize}
    \item \textbf{Spurious probe}: Predicts background (water/land)
    \item \textbf{Core probe}: Predicts bird type (waterbird/landbird)
\end{itemize}

If simplicity bias operates, spurious probe accuracy should exceed core probe accuracy at early epochs.

\subsection{Implementation Details}

\begin{itemize}
    \item \textbf{Model}: ResNet-18 with ImageNet pretrained weights
    \item \textbf{Training}: 100 epochs, SGD (lr=0.001, momentum=0.9, weight decay=$10^{-4}$)
    \item \textbf{Checkpoints}: Epochs 5, 20, 50, 81, 100
    \item \textbf{Seed}: 42
\end{itemize}
